\documentclass[a4paper,10pt,twoside]{article}

\usepackage[top=1in, bottom=1in, left=1in, right=1in]{geometry}
\usepackage[utf8]{inputenc}
\usepackage[spanish,es-ucroman,es-noquoting]{babel}
\usepackage{setspace}
\usepackage{fancyhdr}
\usepackage{lastpage}
\usepackage{amsmath}
\usepackage{amsfonts}
\usepackage{verbatim}
\usepackage{graphicx}
\usepackage{float}
\usepackage{algpseudocode}
\usepackage{enumitem} % Provee macro \setlist
\usepackage[toc, page]{appendix}


%%%%%%%%%% Configuración de Fancyhdr - Inicio %%%%%%%%%%
\pagestyle{fancy}
\thispagestyle{fancy}
\lhead{Trabajo Práctico 2 · Sistemas Operativos}
\rhead{Delgado · Lovisolo · Petaccio}
\renewcommand{\footrulewidth}{0.4pt}
\cfoot{\thepage /\pageref{LastPage}}

\fancypagestyle{caratula} {
   \fancyhf{}
   \cfoot{\thepage /\pageref{LastPage}}
   \renewcommand{\headrulewidth}{0pt}
   \renewcommand{\footrulewidth}{0pt}
}
%%%%%%%%%% Configuración de Fancyhdr - Fin %%%%%%%%%%


%%%%%%%%%% Configuración de Algorithmic - Inicio %%%%%%%%%%
% Entorno propio para customizar la presentación del pseudocódigo
\newenvironment{pseudo}[1][]{%
    \vspace{0.5em}%
    \begin{algorithmic}%
}
{%
    \end{algorithmic}%
    \vspace{0.5em}%
}

% Llamada a función para usar así: \Fn{Foo}{bar, baz}.
% Produce \textsc{Foo}$(bar, baz)$.
\newcommand{\Fn}[2]{\textsc{#1}$(#2)$}

% Cláusula return para usar así: \PReturn foo
\newcommand{\PReturn}[1]{\textbf{return} $#1$}

% Cláusula break
\newcommand{\Break}{\textbf{break}}

% Operadores lógicos
\newcommand{\PAnd}{\textbf{and} }
\newcommand{\POr}{\textbf{or} }
\newcommand{\PNot}{\textbf{not} }

% Booleanos
\newcommand{\PTrue}{\textbf{true} }
\newcommand{\PFalse}{\textbf{false} }

% Null
\newcommand{\Null}{\textbf{null} }

% Conectivo 'in' para usar así: \ForAll{$foo$ \In $bar$}
\newcommand{\In}{\textbf{in} }

% Conectivo 'to' para usar así: \For{$i = 1$ \In $n$}
\newcommand{\To}{\textbf{to} }

% Complejidades
\newcommand{\Ode}[1]{\hfill $O(#1)$}
%%%%%%%%%% Configuración de Algorithmic - Fin %%%%%%%%%%


%%%%%%%%%% Configuración de Appendix - Inicio %%%%%%%%%%
% Asigna la traducción de la palabra 'Appendices'.
\renewcommand{\appendixtocname}{Apéndices}
\renewcommand{\appendixpagename}{Apéndices}
%%%%%%%%%% Configuración de Appendix - Fin %%%%%%%%%%


%%%%%%%%%% Miscelánea - Inicio %%%%%%%%%%
% Evita que el documento se estire verticalmente para ocupar el espacio vacío
% en cada página.
\raggedbottom

% Deshabilita sangría en la primer línea de un párrafo.
\setlength{\parindent}{0em}

% Separación entre párrafos.
\setlength{\parskip}{0.5em}

% Separación entre elementos de listas.
\setlist{itemsep=0.5em}
%%%%%%%%%% Miscelánea - Fin %%%%%%%%%%


\begin{document}


%%%%%%%%%%%%%%%%%%%%%%%%%%%%%%%%%%%%%%%%%%%%%%%%%%%%%%%%%%%%%%%%%%%%%%%%%%%%%%%
%% Carátula                                                                  %%
%%%%%%%%%%%%%%%%%%%%%%%%%%%%%%%%%%%%%%%%%%%%%%%%%%%%%%%%%%%%%%%%%%%%%%%%%%%%%%%


\thispagestyle{caratula}

\begin{center}

\includegraphics[height=2cm]{DC.png} 
\hfill
\includegraphics[height=2cm]{UBA.jpg} 

\vspace{2cm}

Departamento de Computación,\\
Facultad de Ciencias Exactas y Naturales,\\
Universidad de Buenos Aires

\vspace{4cm}

\begin{Huge}
Trabajo Práctico 2
\end{Huge}

\vspace{0.5cm}

\begin{Large}
Sistemas Operativos
\end{Large}

\vspace{1cm}

Segundo Cuatrimestre de 2013

\vspace{4cm}

\begin{tabular}{|c|c|c|}
\hline
Apellido y Nombre & LU & E-mail\\
\hline
Alejandro Nahuel Delgado & 601/11 & nahueldelgado@gmail.com\\
Leandro Lovisolo         & 645/11 & leandro@leandro.me\\
Lautaro José Petaccio    & 443/11 & lausuper@gmail.com\\
\hline
\end{tabular}

\end{center}

\newpage


%%%%%%%%%%%%%%%%%%%%%%%%%%%%%%%%%%%%%%%%%%%%%%%%%%%%%%%%%%%%%%%%%%%%%%%%%%%%%%%
%% Índice                                                                    %%
%%%%%%%%%%%%%%%%%%%%%%%%%%%%%%%%%%%%%%%%%%%%%%%%%%%%%%%%%%%%%%%%%%%%%%%%%%%%%%%


\tableofcontents

\newpage


%%%%%%%%%%%%%%%%%%%%%%%%%%%%%%%%%%%%%%%%%%%%%%%%%%%%%%%%%%%%%%%%%%%%%%%%%%%%%%%
%% Readers-Writers Lock                                                      %%
%%%%%%%%%%%%%%%%%%%%%%%%%%%%%%%%%%%%%%%%%%%%%%%%%%%%%%%%%%%%%%%%%%%%%%%%%%%%%%%


\section{Readers-Writers Lock}

Tomamos la implementación del readers-writers lock libre de inanición de escritura presentada en el libro \textit{The Little Book of Semaphores}, versión 2.1.5, por Allen B. Downey.

Este libro utiliza la analogía de una habitación en la que pueden haber o varios lectores a la vez y ningún escritor, o un único escritor por vez y ningún lector; y un interruptor que prende o apaga la luz de la habitación para señalar si ésta está ocupada o no.

Durante la creación del lock se inicializan las siguientes variables:

\begin{pseudo}
    \State $readers \leftarrow 0$ \Comment{Cantidad de \textbf{lectores} actualmente en la habitación.}
    \State $turnstile \leftarrow$ \Fn{Semaphore}{1} \Comment{Molinete que deja pasar lectores de a uno por vez.}
    \State $room\_empty \leftarrow$ \Fn{Semaphore}{1} \Comment{Vale 1 si la habitación está vacía, o 0 en caso contrario.}
    \State $readers\_mutex \leftarrow$ \Fn{Semaphore}{1} \Comment{Mutex para acceder a la variable $readers$.}
\end{pseudo}

Las operaciones del lock quedan definidas como a continuación.

\begin{pseudo}
    \Procedure{Read-Lock}{}
        \State \Fn{Wait}{turnstile} \Comment{Paso por el molinete, un lector a la vez.}
        \State \Fn{Signal}{turnstile}
        \State
        \State \Fn{Wait}{readers\_mutex}
        \State $readers \leftarrow readers + 1$ \Comment{Aumento en 1 la cantidad de lectores.}
        \If{readers = 1} \Comment{¿Soy el primer lector en llegar?}
            \State \Fn{Wait}{room\_empty} \Comment{Espero a que se vacíe la habitación, luego entro y prendo al luz.}
        \EndIf
        \State \Fn{Signal}{readers\_mutex}
    \EndProcedure
\end{pseudo}

\begin{pseudo}
    \Procedure{Read-Unlock}{}
        \State \Fn{Wait}{readers\_mutex}
        \State $readers \leftarrow readers - 1$ \Comment{Decremento en 1 la cantidad de lectores.}
        \If{readers = 0} \Comment{¿Soy el último lector en irse?}
            \State \Fn{Signal}{room\_empty} \Comment{Apago la luz al salir.}
        \EndIf
        \State \Fn{Signal}{readers\_mutex}
    \EndProcedure
\end{pseudo}

\begin{pseudo}
    \Procedure{Write-Lock}{}
        \State \Fn{Wait}{turnstile} \Comment{Bloqueo el molinete para que no ingresen nuevos lectores.}
        \State \Fn{Wait}{room\_empty} \Comment{Espero a que se vacíe la habitación; luego entro y prendo la luz.}
    \EndProcedure
\end{pseudo}

\begin{pseudo}
    \Procedure{Write-Unlock}{}
        \State \Fn{Signal}{turnstile} \Comment{Permito el ingreso a lectores.}
        \State \Fn{Signal}{room\_empty} \Comment{Apago la luz al salir (señalo que la habitación está vacía.)}
    \EndProcedure
\end{pseudo}

Cuando llega un lector (\Fn{Read-Lock}{}), éste pasa exitosamente por el molinete (\Fn{Wait}{turnstile}) pues el semáforo $turnstile$ vale inicialmente 1 y restaura el valor original del semáforo (\Fn{Signal}{turnstile}) para que subsiguientes lectores puedan ingresar. Luego incrementa en 1 la cantidad de lectores, y si resulta que es el primer lector en ingresar, espera que se vacíe la habitación en caso de estar ocupada por algún escritor (\Fn{Wait}{room\_empty}) y prende la luz al entrar ($room\_empty$ vale 0 luego del \textsc{Wait}.)

Al retirarse (\Fn{Read-Unlock}{}), el lector decrementa en 1 la cantidad de lectores, y en caso de ser el último, apaga la luz al salir (\Fn{Signal}{room\_empty}.)

Cuando llega un escritor (\Fn{Write-Lock}{}), éste bloquea el molinete para que no ingresen nuevos lectores (\Fn{Wait}{turnstile}) y en caso de que haya algún lector en la habitación, espera hasta que se retire el último lector (\Fn{Wait}{room\_empty}) y prende la luz al entrar ($room\_empty$ vale 0 luego del \textsc{Wait}.) Al bloquear el ingreso a nuevos lectores \textbf{nos aseguramos que no ocurra inanición de escritura}, pues de lo contrario podría darse el caso de tener un número muy grande de lectores ingresando constantemente a la habitación, haciendo que ésta nunca quede vacía.

Finalmente, al retirarse un escritor (\Fn{Read-Unlock}{}), éste desbloquea el molinete para permitir el paso a los lectores (\Fn{Signal}{turnstile}) y apaga la luz al salir (\Fn{Signal}{room\_empty}). 

%%%%%%%%%%%%%%%%%%%%%%%%%%%%%%%%%%%%%%%%%%%%%%%%%%%%%%%%%%%%%%%%%%%%%%%%%%%%%%%
%% Server                                                                    %%
%%%%%%%%%%%%%%%%%%%%%%%%%%%%%%%%%%%%%%%%%%%%%%%%%%%%%%%%%%%%%%%%%%%%%%%%%%%%%%%


\section{Server}

Pendiente


%%%%%%%%%%%%%%%%%%%%%%%%%%%%%%%%%%%%%%%%%%%%%%%%%%%%%%%%%%%%%%%%%%%%%%%%%%%%%%%
%% Modelo                                                                    %%
%%%%%%%%%%%%%%%%%%%%%%%%%%%%%%%%%%%%%%%%%%%%%%%%%%%%%%%%%%%%%%%%%%%%%%%%%%%%%%%


\section{Modelo}


\subsection{Modelo::Modelo}

Pendiente


\subsection{Modelo::\char`\~Modelo}

Pendiente


\subsection{Modelo::agregarJugador}

\begin{pseudo}
    \Procedure{Modelo::agregarJugador}{$nombre$}
        \State
        \State Protegemos $jugando$ con un read lock para evitar que comience el juego mientras estamos agregando al nuevo jugador.
        \State
        \State \Fn{rlock}{lock\_jugando}
        \If{$jugando \neq$ \texttt{SETUP}}
            \State \Fn{runlock}{lock\_jugando}
            \State return -\texttt{ERROR\_JUEGO\_EN\_PROGRESO}
        \EndIf
        \State  
        \State $nuevoid \leftarrow$ 0
        \State $agregado \leftarrow$ \PFalse
        \For{$nuevoid \leftarrow$ 0 \To $max\_jugadores$ - 1}
            \State
            \State Protegemos el jugador $nuevoid$ con un write lock para evitar que al agregarse otro jugador en simultáneo, éste no utilice el mismo id.
            \State
            \State \Fn{wlock}{locks\_jugadores[nuevoid]}
            \If{$jugadores[nuevoid] =$ \Null}
                \State $jugadores[nuevoid] \leftarrow$ \Fn{jugador}{nombre}
                \State \Fn{wunlock}{locks\_jugadores[nuevoid]}
                \State agregado $\leftarrow$ \PTrue
                \State \Break
            \EndIf
            \State \Fn{wunlock}{locks\_jugadores[nuevoid]}
        \EndFor
         \If{\PNot $agregado$}
            \State \Fn{runlock}{lock\_jugando}
            \State \PReturn{- \texttt{ERROR\_MAX\_JUGADORES}}
        \EndIf
        \State
        \State \Fn{wlock}{lock\_cantidad\_jugadores}
        \State $cantidad\_jugadores \leftarrow cantidad\_jugadores + 1$
        \State \Fn{wunlock}{lock\_cantidad\_jugadores}
        \State
        \State \Fn{runlock}{lock\_jugando}
        \State \PReturn{nuevoid}
    \EndProcedure
\end{pseudo}


\subsection{Modelo::ubicar}

\begin{pseudo}
    \Procedure{Modelo::ubicar}{$t\_id, xs, ys, tamanio$}
        \State
        \State Protegemos $jugando$ con un read lock para evitar que comience el juego mientras estamos ubicando un barco.
        \State
        \State \Fn{wlock}{lock\_jugando}
        \If{jugando $\neq$ \texttt{SETUP}}
            \State \Fn{wunlock}{lock\_jugando}
            \State \PReturn{- \texttt{ERROR\_JUEGO\_EN\_PROGRESO}}
        \EndIf
        \State
        \State Protegemos el jugador contra lecturas y escrituras simultáneas.
        \State
        \State \Fn{wlock}{locks\_jugadores[t\_id]}
        \If{jugadores[t\_id] = \Null}
            \State \Fn{wunlock}{locks\_jugadores[t\_id]}
            \State \Fn{wunlock}{lock\_jugando}
            \State \PReturn{- \texttt{ERROR\_JUGADOR\_INEXISTENTE}}
        \EndIf
        \State
        \State $retorno \leftarrow jugadores[t\_id].$\Fn{ubicar}{xs, ys, tamanio}
        \State
        \If{$retorno \neq$ \texttt{ERROR\_NO\_ERROR}}
            \State \Fn{borrar\_barcos}{t\_id}
        \EndIf
        \State
        \If{$retorno = $ \texttt{ERROR\_NO\_ERROR} \PAnd $jugadores[t\_id].$\Fn{listo}{}} \Comment{Si el jugador está listo}
            \State \Fn{wlock}{lock\_jugadores\_listos}
            \State $jugadores\_listos \leftarrow jugadores\_listos + 1$
            \State
            \If{$jugadores\_listos = max\_jugadores$} \Comment{Si ya están listos todos los jugadores}
                \State \Fn{\_empezar}{}
            \EndIf
            \State \Fn{wunlock}{lock\_jugadores\_listos}
        \EndIf
        \State
        \State \Fn{wunlock}{locks\_jugadores[t\_id]}
        \State \Fn{wunlock}{locks\_jugando}
        \State \PReturn{retorno}
    \EndProcedure
\end{pseudo}


\subsection{Modelo::borrar\_barcos}

Pendiente


\subsection{Modelo::empezar}

Separamos este método en dos partes: \textsc{Modelo::empezar}, que toma write lock sobre $jugando$ y llama a \textsc{Modelo::\_empezar}, que asume que ya se tiene el write lock sobre $jugando$ y realiza la operación de empezar. El método \textsc{Modelo::empezar} se llama desde la clase \textsc{Decodificador}, que no tiene acceso a dicho lock directamente, mientras que \textsc{Modelo::\_empezar} se llama desde \textsc{Modelo::ubicar}, que ya toma lock de escritura sobre $jugando$ para proteger operaciones anteriores a \textsc{Modelo::\_empezar}.

Si no se abriera este método en los dos casos mencionados, tendríamos un único método que toma write lock sobre $jugando$, pues como \textsc{Decodificador} no tiene acceso a ese lock, no queda otra opción. Entonces \textsc{Modelo::ubicar} produciría dos llamadas consecutivas \Fn{wlock}{lock\_jugando}, generando así un deadlock. Separando este método en dos casos distintos evitamos el deadlock.

\begin{pseudo}
    \Procedure{Modelo::empezar}{}
        \State \Fn{wlock}{lock\_jugando}
        \State $out \leftarrow$ \Fn{\_empezar}{}
        \State \Fn{wunlock}{lock\_jugando}
        \State \PReturn{out}
    \EndProcedure
\end{pseudo}

\begin{pseudo}
    \Procedure{Modelo::\_empezar}{}
    	\If{$jugando \neq$ \texttt{SETUP}}
    		\State \PReturn{-\texttt{ERROR\_JUEGO\_EN\_PROGRESO}}
    	\EndIf
    	\For{$i \leftarrow 0$ \To $max\_jugadores - 1$}
    		\If{$jugadores[i] \neq$ \Null}
    			\State $evento \leftarrow$ \Fn{evento}{0, i, 0, 0, \texttt{EVENTO\_START}}
                \State
                \State Tomamos write lock de la cola de eventos del jugador para evitar que ésta se lea desde el ciclo principal de su thread en server.cpp mientras agregamos el nuevo evento.
                \State
    			\State \Fn{wlock}{locks\_eventos[i]}
    			\State $eventos[i].$\Fn{push}{evento}
    			\State \Fn{wunlock}{locks\_eventos[i]}
    		\EndIf
    	\EndFor
    	\State $jugando \leftarrow$ \texttt{DISPAROS}
    	\State \PReturn{\texttt{ERROR\_NO\_ERROR}}
    \EndProcedure
\end{pseudo}


\subsection{Modelo::finalizar}

\begin{pseudo}
    \Procedure{modelo::finalizar}{}
    	\If{$jugando \neq$ \texttt{DISPAROS}}
            \State \PReturn{-\texttt{ERROR\_JUEGO\_NO\_COMENZADO}}
        \EndIf
    	\For{$i \leftarrow 0$ \To $max\_jugadores - 1$}
    		\If{$jugadores[i] \neq$ \Null}
    			\State $evento  \leftarrow$ \Fn{evento}{0, i, 0, 0, \texttt{EVENTO\_START}}
                \State
                \State Tomamos write lock de la cola de eventos del jugador para evitar que ésta se lea desde el ciclo principal de su thread en server.cpp mientras agregamos el nuevo evento.
                \State
    			\State \Fn{wlock}{locks\_eventos[i]}
    			\State $eventos[i].$\Fn{push}{evento}
    			\State \Fn{wunlock}{locks\_eventos[i]}
    		\EndIf
    	\EndFor
    	\State $jugando \leftarrow$ \texttt{FINALIZADO}
    	\State \PReturn{\texttt{ERROR\_NO\_ERROR}}
    \EndProcedure
\end{pseudo}


\subsection{Modelo::ack}

\begin{pseudo}
    \Procedure{modelo::ack}{}
        \State
        \State Tomamos write lock para evitar lecturas mientras registramos el ack.
        \State
    	\State \Fn{wlock}{locks\_jugadores[s_id]}
    	\State $retorno \leftarrow$ jugadores[s\_id].\Fn{ack}{}
    	\State \Fn{wunlock}{locks\_jugadores[s\_id]}
    	\State \PReturn{retorno}
    \EndProcedure
\end{pseudo}


\subsection{Modelo::termino}

\begin{pseudo}
    \Procedure{termino}{}
    	\State \Fn{rlock}{lock\_jugando}
    	\If{$jugando$ = \texttt{SETUP}}
    		\State \Fn{runlock}{lock\_jugando}
    		\State \PReturn{\PFalse}
    	\EndIf
    	\State \Fn{runlock}{lock\_jugando}
        \State
        \For{$i \leftarrow 0$ \To $max\_jugadores$ - 1}
        	\State \Fn{rlock}{locks\_jugadores[i]}
            \If{\PNot $jugadores[i].$\Fn{termino}{}}
            	\State \Fn{runlock}{locks\_jugadores[i]}
                \State \PReturn{\PFalse}
            \EndIf
            \State \Fn{runlock}{locks\_jugadores[i]}
        \EndFor
        \State \PReturn{\PTrue}
    \EndProcedure
\end{pseudo}


\subsection{Modelo::quitarJugador}

\begin{pseudo}
    \Procedure{Modelo::quitarJugador}{$s\_id$}
    	\State \Fn{rlock}{lock\_jugando}
    	\If{$jugando \neq$ \texttt{SETUP}}
    		\State \Fn{runlock}{lock\_jugando}
    		\State \PReturn{-\texttt{ERROR\_JUEGO\_EN\_PROGRESO}}
    	\EndIf
        \State
    	\State \Fn{wlock}{locks\_jugadores[s\_id]}
    	\If{$jugadores[s\_id]$ = \Null}
    		\State \Fn{wunlock}{locks\_jugadores[s\_id]}
    		\State \Fn{runlock}{lock\_jugando}
    		\State \PReturn{-\texttt{ERROR\_JUGADOR\_INEXISTENTE}}
    	\EndIf
        \State
    	\State \Fn{delete}{jugadores[s\_id]}
        \State $jugadores[s\_id] \leftarrow$ \Null
        \State
    	\State \Fn{wunlock}{locks\_jugadores[s\_id]}
    	\State \Fn{runlock}{lock\_jugando}
    	\State \PReturn{\texttt{ERROR\_NO\_ERROR}}
    \EndProcedure
\end{pseudo}


\subsection{Modelo::apuntar}

\begin{pseudo}
    \Procedure{apuntar}{$s\_id, t\_id, x, y, *eta$}
        \State \Fn{rlock}{lock\_jugando}
        \If{$jugando \neq$ \texttt{DISPAROS}}
            \State \Fn{runlock}{lock\_jugando}
            \State \PReturn{-\texttt{ERROR\_JUEGO\_NO\_COMENZADO}}
        \EndIf
        \State
        \State \Fn{wlock}{locks\_jugadores[s\_id]}
        \State \Fn{wlock}{locks\_jugadores[t\_id]}
        \State
        \If{$jugadores[s\_id] = $ \Null \POr $jugadores[t\_id] = $ \Null}
            \State \Fn{runlock}{lock\_jugando}
            \State \Fn{wunlock}{locks\_jugadores[s\_id]}
            \State \Fn{wunlock}{locks\_jugadores[t\_id]}            
            \State \PReturn{-\texttt{ERROR\_JUGADOR\_INEXISTENTE}}
        \EndIf
        \State
        \If{$jugadores[s\_id].$\Fn{esta\_vivo}{}}
            \State \Fn{runlock}{lock\_jugando}
            \State \Fn{wunlock}{locks\_jugadores[s\_id]}
            \State \Fn{wunlock}{locks\_jugadores[t\_id]}            
            \State \PReturn{-\texttt{ERROR\_JUGADOR\_HUNDIDO}}        
        \EndIf
        \State
        \State $retorno \leftarrow$ \texttt{RESULTADO\_APUNTADO\_DENEGADO}
        \State
        \State \Fn{wlock}{locks\_tiros[s\_id]}
        \If{$tiros[s\_id].$\Fn{es\_posible\_apuntar}{}}
            \State $retorno \leftarrow jugadores[t\_id].$\Fn{apuntar}{s\_id, x, y}
            \If{$retorno = $ \texttt{RESULTADO\_APUNTADO\_ACEPTADO}}
                \State $*eta \leftarrow tiros[s\_id].$\Fn{tirar}{t\_id, x, y}
                \State $nuevoevento \leftarrow $ \Fn{evento}{s\_id, t\_id, x, y, \texttt{CASILLA\_EVENTO\_INCOMING}}
                \State
                \State \Fn{wlock}{locks\_eventos[t\_id]}
                \State $eventos[t\_id].$\Fn{push}{nuevoevento}
                \State \Fn{wunlock}{locks\_eventos[t\_id]}
            \EndIf
        \EndIf
        \State \Fn{wunlock}{locks\_tiros[s\_id]}
        \State
        \State \Fn{wunlock}{locks\_jugadores[s\_id]}
        \State \Fn{wunlock}{locks\_jugadores[t\_id]}
        \State \Fn{runlock}{lock\_jugando}
        \State \PReturn{retorno}
    \EndProcedure
\end{pseudo}


\subsection{Modelo::tocar}

\begin{pseudo}
    \Procedure{tocar}{$s\_id, t\_id$}
        \State \Fn{wlock}{lock\_jugando}
        \If{$jugando \neq $ \texttt{DISPAROS}}
            \State \Fn{wunlock}{lock\_jugando}
            \State \PReturn{- \texttt{ERROR\_JUEGO\_NO\_COMENZADO}}
        \EndIf
        \State
        \If{$jugadores[s\_id] =$ \Null \POr $jugadores[t\_id] = $ \Null}
            \State \Fn{wunlock}{lock\_jugando}
            \State \PReturn{- \texttt{ERROR\_JUGADOR\_INEXISTENTE}}
        \EndIf
        \State
        \State $retorno \leftarrow$ -\texttt{ERROR\_ETA\_NO\_TRANSCURRIDO}
        \State
        \State \Fn{wlock}{locks\_tiros[s\_id]}
        \State
        \If{$tiros[s\_id].$\Fn{es\_posible\_tocar}{}}
            \State \Fn{wlock}{locks\_jugadores[t\_id]}
            \State \Fn{wlock}{locks\_jugadores[s\_id]}
            \State
            \State $x \leftarrow tiros[s\_id].x$
            \State $y \leftarrow tiros[s\_id].y$
            \State $retorno, murio \leftarrow jugadores[t\_id].$\Fn{tocar}{s\_id, x, y} \Comment{Devuelve un código de error y un}
            \State \hfill booleano que indica si murió o no.
            \State
            \If{$retorno = $ \texttt{EMBARCACION\_RESULTADO\_TOCADO} \POr \\
                \hspace{3.8em} $retorno = $ \texttt{EMBARCACION\_RESULTADO\_HUNDIDO} \POr \\
                \hspace{3.8em} $retorno = $ \texttt{EMBARCACION\_RESULTADO\_HUNDIDO\_M} \POr \\
                \hspace{3.8em} $retorno = $ \texttt{EMBARCACION\_RESULTADO\_AGUA} \POr \\
                \hspace{3.8em} $retorno = $ \texttt{EMBARCACION\_RESULTADO\_AGUA\_H}}
                \State
                \State $tiros[s\_id].estado \leftarrow$ \texttt{TIRO\_LIBRE}
                \State $evento \leftarrow$ \Fn{evento}{s\_id, t\_id, x, y, retorno}
                \State
                \State \Fn{wlock}{locks\_eventos[t\_id]} \Comment{Evento para el tirado}
                \State $eventos[t\_id].$\Fn{push}{evento}
                \State \Fn{wunlock}{locks\_eventos[t\_id]}
                \State
                \State \Fn{wlock}{locks\_eventos[s\_id]} \Comment{Evento para el tirador}
                \State $eventos[s\_id].$\Fn{push}{evento}
                \State \Fn{wunlock}{locks\_eventos[s\_id]}
                \State
                \If{$murio$}
                    \State \Fn{wlock}{lock\_cantidad\_jugadores}
                    \State $cantidad\_jugadores \leftarrow cantidad\_jugadores - 1$
                    \If{$cantidad\_jugadores = 1$}
                        \State \Fn{finalizar}{}
                    \EndIf
                    \State \Fn{wunlock}{lock\_cantidad\_jugadores}
                \EndIf
            \EndIf
            \If{$retorno =$ \texttt{EMBARCACION\_RESULTADO\_HUNDIDO}}
                \State $jugadores[s\_id].$\Fn{agregar\_puntaje}{\texttt{PUNTAJE\_HUNDIDO}}
            \ElsIf{$retorno =$ \texttt{EMBARCACION\_RESULTADO\_HUNDIDO\_M}}
                \State $jugadores[s\_id].$\Fn{agregar\_puntaje}{\texttt{PUNTAJE\_HUNDIDO} + \texttt{PUNTAJE\_MISMO\_JUGADOR}}
            \ElsIf{$retorno =$ \texttt{EMBARCACION\_RESULTADO\_TOCADO}}
                \State $jugadores[s\_id].$\Fn{agregar\_puntaje}{\texttt{PUNTAJE\_TOCADO}}
            \ElsIf{$retorno =$ \texttt{EMBARCACION\_RESULTADO\_AGUA\_H}}
                \State $jugadores[s\_id].$\Fn{agregar\_puntaje}{\texttt{PUNTAJE\_MAGALLANES}}
            \EndIf
            \State
            \State \Fn{wunlock}{locks\_jugadores[t\_id]}
            \State \Fn{wunlock}{locks\_jugadores[s\_id]}
        \EndIf
        \State
        \State \Fn{wunlock}{locks\_tiros[s\_id]}
        \State \Fn{wunlock}{lock\_jugando}
        \State
        \State \PReturn{retorno}
    \EndProcedure
\end{pseudo}


\subsection{Modelo::hayEventos}

\begin{pseudo}
    \Procedure{hayEventos}{$s\_id$}
        \State \Fn{rlock}{locks\_eventos[s\_id]}
        \State $eventos \leftarrow eventos[s\_id]$.\Fn{size}{} $> 0$
        \State \Fn{runlock}{locks\_eventos[s\_id]}
        \State \PReturn{eventos}
    \EndProcedure
\end{pseudo}


\subsection{Modelo::dameEvento}

\begin{pseudo}
    \Procedure{dameEvento}{$s\_id$}
        \State \Fn{wlock}{locks\_eventos[s\_id]}
        \State $retorno \leftarrow eventos[s\_id]$.\Fn{front}{}
        \State $eventos[s_id]$.\Fn{pop}{}
        \State \Fn{wunlock}{locks\_eventos[s\_id]}
        \State \PReturn{eretorno}
    \EndProcedure
\end{pseudo}


\subsection{Modelo::actualizar\_jugador}

\begin{pseudo}
    \Procedure{actualizar\_jugador}{$s\_id$}
        \State $tocado \leftarrow$ \Fn{tocar}{s\_id, tiros[s\_id].t\_id}
        \State \Fn{wlock}{locks\_eventos[s\_id]}
        \If{$eventos[s\_id]$.\Fn{size}{} $> 0$}
            \State $retorno \leftarrow eventos[s\_id]$.\Fn{front}{}
            \State $eventos[s\_id]$.\Fn{pop}{}
            \State \Fn{wunlock}{locks\_eventos[s\_id]}
            \State \PReturn{retorno}
        \Else
            \State $retorno \leftarrow$ \Fn{evento}{s\_id, -1, 0, 0, -tocado}
            \State \Fn{wunlock}{locks\_eventos[s\_id]}
            \State \PReturn{retorno}
        \EndIf
    \EndProcedure
\end{pseudo}


%%%%%%%%%%%%%%%%%%%%%%%%%%%%%%%%%%%%%%%%%%%%%%%%%%%%%%%%%%%%%%%%%%%%%%%%%%%%%%%
%% Jsonificador                                                              %%
%%%%%%%%%%%%%%%%%%%%%%%%%%%%%%%%%%%%%%%%%%%%%%%%%%%%%%%%%%%%%%%%%%%%%%%%%%%%%%%


\section{Jsonificador}


\subsection{Jsonificador::start}

Utilizamos un read lock para que otras operaciones sobre los jugadores no modifiquen al jugador del que estoy
leyendo su información.

\begin{pseudo}
    \For{$i \leftarrow 0$ \To $max\_jugadores$}
       \State \Fn{rlock}{modelo.locks\_jugadores[i]}
       \State Lectura de datos del jugador
       \State \Fn{runlock}{modelo.locks\_jugadores[i]}

    \EndFor
\end{pseudo}

\subsection{Jsonificador::player\_info}

Utilizamos un read lock para que ninguna otra operación escriba mientras leemos los datos del jugador.

\begin{pseudo}
    \State \Fn{rlock}{modelo.locks\_jugadores[id]}
    \State Lectura de datos del jugador
    \State \Fn{runlock}{modelo.locks\_jugadores[id]}
\end{pseudo}


\subsection{Jsonificador::scores}

Protegemos con un read lock al jugador a leer para que ninguna otra operación escriba
sobre éste mientras obtenemos información de este.

\begin{pseudo}
    \For{$i \leftarrow 0$ To $max\_jugadores$}
       \State \Fn{rlock}{modelo->locks\_jugadores[i]}
       \State Utilización de datos del jugador
       \State \Fn{runlock}{modelo->locks\_jugadores[i]}
    \EndFor
\end{pseudo}
\end{document}